\chapter{The Peripheral HAL}
\label{ch:hal}

The hand-written drivers live under \texttt{libs/esp32s3\_hal} as children of a
single root package, \texttt{ESP32S3}: \texttt{ESP32S3.GPIO},
\texttt{ESP32S3.SPI}, and so on. A consumer adds one line to its project file,

\begin{lstlisting}
with "esp32s3_hal.gpr";   --  or a relative path in-repo
\end{lstlisting}

and then \texttt{with}s the drivers it needs. Every driver in this book has been
exercised on real silicon.

\section{Conventions used by every driver}

\begin{description}[style=nextline]
\item[Valid-pin types.] A pin is named with \texttt{ESP32S3.GPIO.Pin\_Id}, a
  subtype whose \emph{static predicate} admits only the pads the ESP32-S3 exposes
  (0--21 and 38--48); the flash/PSRAM pads are excluded. A bad constant pin is a
  compile-time error.
\item[Errors.] Simple drivers report failure through a \texttt{Status} enum or an
  \texttt{out Boolean}; the filesystem and richer layers raise exceptions. Either
  way, failure is explicit.
\item[Task-safety, two shapes.] A peripheral that is a single shared controller
  (SPI host, I2C host, UART port, I2S port, the TWAI/LCD controller) is mediated
  by a protected object plus a limited, non-copyable \texttt{Session} handle: you
  \texttt{Acquire} it, use it, and it is released automatically when it leaves
  scope. A peripheral that is a \emph{pool} of identical units (GDMA channels,
  LEDC/RMT/PCNT/SDM channels, timers, ADC units, MCPWM channels) hands out a
  limited, controlled handle you \texttt{Claim}; ownership makes per-handle
  operations safe without a global lock, and the unit is reclaimed on scope exit.
\item[RAII.] Those \texttt{Session}/\texttt{Channel}/\texttt{Reader} handles are
  controlled types -- they release on scope exit, including during exception
  unwinding -- so they require the \texttt{embedded} or \texttt{full} profile.
  The finalization-free drivers (GPIO, RNG, Temperature, RTC, RTC\_IO, Touch, SHA,
  AES, SDMMC) use no controlled types -- a few serialise internally with a protected
  object instead -- so they work under every profile.
\item[Init-before-use guard.] For the drivers that configure a controller once
  at startup and only later hand out ownership (SPI, I2C, UART, I2S, TWAI, LCD,
  and MCPWM), \texttt{Acquire}/\texttt{Claim} raises \texttt{Not\_Initialized} if
  the port/host/unit was never \texttt{Setup}. Configuration must precede
  ownership; the contract is enforced, not merely documented.
\item[Ownership gateway.] Across all six \texttt{Session} drivers (SPI, I2C,
  UART, I2S, TWAI, LCD), all register access lives in a private \texttt{Engine}
  child, and the configured handle it returns is hidden in a nested package body.
  The only way to reach it is one ownership-checked gateway, \texttt{Owned (S)},
  which raises \texttt{Not\_Owned} unless the caller holds the port. Every
  operation that touches a live controller -- transfers \emph{and} every
  post-\texttt{Setup} reconfiguration (pins, baud, inversion, loopback) -- goes
  through it. So a transfer or a setting change without ownership fails loudly,
  and a new operation cannot be added that skips the check: there is no other
  door to the hardware. This is a project rule for any exclusive driver that can
  be reconfigured after setup (see the UART section for the worked pattern).
\end{description}

\paragraph{Recurring parameter types.} Each driver below lists its parameters
with their allowed values and defaults. A few types recur throughout and are
described once here rather than on every driver:
\begin{description}[style=nextline]
\item[\texttt{Pin\_Id}] a usable GPIO pad -- \texttt{0\,..\,21} or
  \texttt{38\,..\,48} (the gaps are flash/PSRAM/absent pads). A bad constant is a
  compile error.
\item[\texttt{Optional\_Pin}] a \texttt{Pin\_Id} \emph{or} \texttt{No\_Pin} (the
  value $-1$) to leave that line unrouted.
\item[Ownership handles (\texttt{Session}, \texttt{Channel}, \texttt{Reader},
  \texttt{Unit}, \texttt{Timer})] the limited, controlled handles described above:
  you declare one, \texttt{Acquire}/\texttt{Claim} it, pass it to each operation,
  and it releases on scope exit.
\item[\texttt{...\_Hz}, \texttt{Freq}, \texttt{Baud}, \texttt{Bit\_Rate}]
  frequencies / rates in hertz (or bits per second); \texttt{Positive} unless a
  narrower subtype is noted.
\item[\texttt{...\_Percent}] a \texttt{Float} in \texttt{0.0\,..\,100.0}.
\item[Buffers] a transfer takes a buffer \texttt{'Address} and a byte
  \texttt{Length}; the DMA-backed drivers cap a single transfer at \texttt{4095}
  bytes.
\end{description}

\section{GPIO -- digital input and output}\index{GPIO driver}

\texttt{ESP32S3.GPIO} configures a pad's direction, pull and drive strength and
then drives or reads it. \texttt{Set}/\texttt{Clear}/\texttt{Write}/\texttt{Read}
are hardware-atomic (lock-free); \texttt{Configure}/\texttt{Toggle} are serialised
by a protected object.

\begin{lstlisting}
procedure Configure (Pin : Pin_Id; Mode : Pin_Mode;
                     Pull  : Pull_Mode      := Floating;
                     Drive : Drive_Strength := Drive_Medium);
procedure Set    (Pin : Pin_Id);   procedure Clear (Pin : Pin_Id);
procedure Toggle (Pin : Pin_Id);   procedure Write (Pin : Pin_Id; On : Boolean);
function  Read   (Pin : Pin_Id) return Boolean;
\end{lstlisting}

\textbf{Parameters.}
\begin{description}[style=nextline]
\item[\texttt{Pin}] the pad to act on (\texttt{Pin\_Id}).
\item[\texttt{Mode}] direction: \texttt{Input} or \texttt{Output}.
\item[\texttt{Pull}] internal resistor: \texttt{Floating} (default), \texttt{Pull\_Up}
  or \texttt{Pull\_Down}.
\item[\texttt{Drive}] output current: \texttt{Drive\_Weak}, \texttt{Drive\_Medium}
  (default), \texttt{Drive\_Strong}, \texttt{Drive\_Strongest} (roughly
  5 / 10 / 20 / 40\,mA).
\item[\texttt{On}] (\texttt{Write}) the level to drive -- \texttt{True} = high,
  \texttt{False} = low.
\end{description}

\texttt{Configure} always sets the pad up as a software-controlled GPIO (routed
through the GPIO matrix). Handing a pad to a \emph{peripheral} instead is the job
of that peripheral's own \texttt{Configure\_Pins}, which programs the matrix
directly -- the ESP32-S3 GPIO matrix is a full crossbar, so almost any signal
can reach almost any pad.

\textbf{Example 1 -- blink an LED.}
\begin{lstlisting}
with Ada.Real_Time; use Ada.Real_Time;
with ESP32S3.GPIO;  use ESP32S3.GPIO;
procedure Blink is
   Led : constant Pin_Id := 2;
begin
   Configure (Led, Mode => Output);
   loop
      Toggle (Led);
      delay until Clock + Milliseconds (250);
   end loop;
end Blink;
\end{lstlisting}

\textbf{Example 2 -- a button drives an LED.}
\begin{lstlisting}
with ESP32S3.GPIO; use ESP32S3.GPIO;
procedure Button is
   Button_Pin : constant Pin_Id := 9;    --  active-low, internal pull-up
   Led        : constant Pin_Id := 2;
begin
   Configure (Button_Pin, Mode => Input,  Pull => Pull_Up);
   Configure (Led,        Mode => Output);
   loop
      Write (Led, not Read (Button_Pin)); --  pressed (low) -> LED on
   end loop;
end Button;
\end{lstlisting}

\section{RNG -- hardware random numbers}\index{RNG driver}

\begin{lstlisting}
function Read return Word;                 --  a 32-bit random word
procedure Fill (Buffer : out Byte_Array);  --  fill any length
\end{lstlisting}

\textbf{Parameters.} \texttt{Read} takes none and returns a \texttt{Word} (a 32-bit
\texttt{Interfaces.Unsigned\_32}). \texttt{Fill}'s \texttt{Buffer} is an
\texttt{out Byte\_Array} of \emph{any} length -- the call fills every element with
random bytes.

\textbf{Example 1 -- a random word.}
\begin{lstlisting}
with Ada.Text_IO; use Ada.Text_IO;
with ESP32S3.RNG;
procedure RNG_Word is
   X : constant ESP32S3.RNG.Word := ESP32S3.RNG.Read;
begin
   Put_Line (ESP32S3.RNG.Word'Image (X));
end RNG_Word;
\end{lstlisting}

\textbf{Example 2 -- a random nonce.}
\begin{lstlisting}
with ESP32S3.RNG;
procedure RNG_Nonce is
   Nonce : ESP32S3.RNG.Byte_Array (0 .. 15);
begin
   ESP32S3.RNG.Fill (Nonce);   --  16 random bytes
end RNG_Nonce;
\end{lstlisting}

(For true entropy the radio/ADC clocks should be running; see the package spec.)

\section{Temperature -- the on-die sensor}\index{Temperature driver}

\begin{lstlisting}
procedure Initialize (Span : Measure_Range := Range_Minus10_80);
function Read_Celsius return Integer;        --  whole degrees
function Read_Centi_Celsius return Integer;  --  hundredths of a degree
\end{lstlisting}

\textbf{Parameters.} \texttt{Initialize}'s \texttt{Span} (a \texttt{Measure\_Range})
picks the calibrated window -- accuracy is best when it brackets the real
temperature. The five values are \texttt{Range\_Minus10\_80} (default, $-10\,..\,80$\,\textdegree C),
\texttt{Range\_20\_100}, \texttt{Range\_50\_125}, \texttt{Range\_Minus30\_50}
($-30\,..\,50$) and \texttt{Range\_Minus40\_20} ($-40\,..\,20$). The two
\texttt{Read\_*} functions take no parameters; one returns whole degrees, the
other hundredths.

\textbf{Example 1 -- read degrees.}
\begin{lstlisting}
with Ada.Text_IO;        use Ada.Text_IO;
with ESP32S3.Temperature; use ESP32S3.Temperature;
procedure Temp_Degrees is
begin
   Initialize;
   Put_Line ("die temp =" & Integer'Image (Read_Celsius) & " C");
end Temp_Degrees;
\end{lstlisting}

\textbf{Example 2 -- higher resolution and a different span.}
\begin{lstlisting}
with Ada.Text_IO;        use Ada.Text_IO;
with ESP32S3.Temperature; use ESP32S3.Temperature;
procedure Temp_Centi is
begin
   Initialize (Span => Range_50_125);          --  a hotter measurement range
   declare
      Centi : constant Integer := Read_Centi_Celsius;
   begin
      Put_Line ("temp =" & Integer'Image (Centi / 100) & "." &
                Integer'Image (Centi mod 100) & " C");
   end;
end Temp_Centi;
\end{lstlisting}

\section{SPI -- full-duplex master}\index{SPI driver}

\texttt{ESP32S3.SPI} drives either general-purpose host (\texttt{SPI2}/\texttt{SPI3})
over GDMA. Bring the host up and route its shared wires once, then \texttt{Acquire}
a \texttt{Session} per transaction -- which is where each device's own SPI mode and
bit clock are applied, so two devices on one host can differ. \texttt{Transfer}
shifts \texttt{Tx} out while capturing \texttt{Rx}.

\begin{lstlisting}
procedure Setup (Host : SPI_Host);                          --  controller + GDMA
procedure Configure_Pins (Host : SPI_Host; Sclk, Mosi, Miso : Optional_Pin;
                          Cs : Optional_Pin := No_Pin);      --  shared wiring
procedure Acquire  (S : in out Session; Host : SPI_Host;
                    Mode      : SPI_Mode := 0;               --  per-device
                    Clock_Hz  : Positive := 1_000_000;       --  per-device
                    Sclk, Mosi, Miso : Optional_Pin := No_Pin; --  pin override
                    CS_Pin    : Optional_Pin := No_Pin;      --  built-in sw select
                    Select_CB : CS_Select    := null;        --  or app chip select
                    Ctx : System.Address := System.Null_Address);
procedure Select_Device (S : in out Session; On : Boolean);
procedure Transfer (S : Session; Tx, Rx : System.Address; Length : Natural);
procedure Release  (S : in out Session);
\end{lstlisting}

The three host calls split by lifetime and scope, the same shape as the other
buses: \texttt{Setup} (once, controller + GDMA), \texttt{Configure\_Pins} (once,
the shared \texttt{Sclk}/\texttt{Mosi}/\texttt{Miso} every device uses), and
\texttt{Acquire} (per transaction, this device's mode/clock/select). SPI's
\texttt{Setup} carries no mode or clock -- unlike, say, \texttt{I2C.Setup} --
precisely because on SPI those are per-\emph{device}, not per-host.

\textbf{Parameters.}
\begin{description}[style=nextline]
\item[\texttt{Host}] which general-purpose host: \texttt{SPI2} or \texttt{SPI3}.
\item[\texttt{Sclk}, \texttt{Mosi}, \texttt{Miso}, \texttt{Cs}] (\texttt{Configure\_Pins})
  the shared bus lines as \texttt{Optional\_Pin}; \texttt{Cs} is the host's single
  hardware chip-select and defaults to \texttt{No\_Pin} -- with more than one device
  on the bus, give each its own select via \texttt{CS\_Pin} at \texttt{Acquire}.
\item[\texttt{Mode}] (\texttt{Acquire}) this device's SPI clock mode \texttt{0\,..\,3}
  (the CPOL/CPHA combinations); default \texttt{0}, applied under the hold.
\item[\texttt{Clock\_Hz}] (\texttt{Acquire}) this device's bit clock in Hz; default
  \texttt{1\_000\_000}, clamped to roughly 80\,kHz\,..\,80\,MHz.
\item[\texttt{Sclk}, \texttt{Mosi}, \texttt{Miso}] (\texttt{Acquire}) usually
  \texttt{No\_Pin} = keep the \texttt{Configure\_Pins} routing; set them only for a
  device wired to a \emph{different} set of pads on the same controller (rare).
\item[\texttt{CS\_Pin}] a built-in software chip-select: name one GPIO and the
  driver owns it, driving it active-low and holding it asserted across the whole
  transaction. The common case for a shared bus; defaults to \texttt{No\_Pin}.
\item[\texttt{Select\_CB}, \texttt{Ctx}] an application chip-select callback and its
  opaque context, for a select that is \emph{not} one plain GPIO (a decoder, an
  I/O-expander line); leave all three out for the hardware \texttt{CS0}.
\item[\texttt{Tx}, \texttt{Rx}, \texttt{Length}] transmit/receive buffer
  \texttt{'Address}es and a byte count, \texttt{1\,..\,4095}; full-duplex, so both
  move \texttt{Length} bytes together.
\item[\texttt{Pad}] (\texttt{Enable\_Loopback}) one pad that folds MOSI back to MISO
  for the wiring-free self-test.
\end{description}

\textbf{Example 1 -- write a command to a device.}
\begin{lstlisting}
with Interfaces;
with ESP32S3.SPI; use ESP32S3.SPI;
procedure SPI_Write is
   S  : Session;
   Tx : aliased array (0 .. 2) of Interfaces.Unsigned_8 := (16#9F#, 0, 0);
   Rx : aliased array (0 .. 2) of Interfaces.Unsigned_8;
begin
   Setup (SPI2);
   Configure_Pins (SPI2, Sclk => 12, Mosi => 11, Miso => 13, Cs => 10);
   Acquire (S, SPI2, Mode => 0, Clock_Hz => 8_000_000);  --  this device's mode/clock
   Transfer (S, Tx'Address, Rx'Address, 3);   --  Rx now holds the response
end SPI_Write;                                 --  S releases the host here
\end{lstlisting}

\textbf{Example 2 -- wiring-free self-test via internal loopback.}
\begin{lstlisting}
with Interfaces;
with ESP32S3.SPI; use ESP32S3.SPI;
procedure SPI_Loopback is
   S  : Session;
   Tx : aliased array (0 .. 31) of Interfaces.Unsigned_8 := (others => 16#A5#);
   Rx : aliased array (0 .. 31) of Interfaces.Unsigned_8 := (others => 0);
begin
   Setup (SPI2);
   Enable_Loopback (SPI2, Pad => 5);   --  MOSI->MISO through one pad
   Acquire (S, SPI2);
   Transfer (S, Tx'Address, Rx'Address, 32);
   --  Rx = Tx
end SPI_Loopback;
\end{lstlisting}

\textbf{Per-device chip select.}\index{chip-select}\index{shared SPI bus}
The host has one hardware chip-select (\texttt{CS0}); to put several devices on one
bus, each device brings its \emph{own} select instead. The simple case is one
GPIO: pass its pin as \texttt{CS\_Pin} to \texttt{Acquire} and the driver owns it
-- it suppresses the hardware \texttt{CS0} for that session, parks the GPIO high,
and \texttt{Select\_Device} drives it low to assert and high to release. Crucially
the line is toggled only at \texttt{Select\_Device}, never per \texttt{Transfer}, so
it stays asserted across a whole multi-transfer transaction (a streamed flash read
keeps \texttt{CS} low the entire time) -- which the hardware \texttt{CS0}, pulsed
once per transfer, cannot do. This is how the W25Q flash shares the SPI2 bus with
the W5500 in Chapter~\ref{ch:storage}.

\begin{lstlisting}
   Setup (SPI2);
   Configure_Pins (SPI2, Sclk => 12, Mosi => 11, Miso => 13);  --  no hardware Cs
   Acquire (S, SPI2, Clock_Hz => 8_000_000, CS_Pin => 10);  --  drive IO10 active-low
   Select_Device (S, On => True);      --  assert IO10
   Transfer (S, Tx'Address, Rx'Address, 2);
   Select_Device (S, On => False);     --  deassert IO10
\end{lstlisting}

\textbf{When the select is not one GPIO.} A select behind a 3:8 decoder or an
I/O-expander cannot be expressed as a single pin. For those, leave
\texttt{CS\_Pin} at \texttt{No\_Pin} and pass a \texttt{CS\_Select} callback (with
an opaque \texttt{Ctx}) instead; \texttt{Select\_Device} then calls back to assert
and deassert, with exactly the same held-across-the-transaction semantics. The
callback must be \emph{library-level and closure-free} (this HAL builds under
\texttt{No\_Implicit\_Dynamic\_Code}, so a nested closure would emit a faulting
trampoline; per-device state travels in \texttt{Ctx}).

\begin{lstlisting}
--  Library-level, closure-free: drive IO10 low to select, high to release.
procedure CS_IO10 (Ctx : System.Address; Active : Boolean) is
begin
   if Active then ESP32S3.GPIO.Clear (10); else ESP32S3.GPIO.Set (10); end if;
end CS_IO10;

   Acquire (S, SPI2, Select_CB => CS_IO10'Access);   --  bring our own select
\end{lstlisting}

\textbf{Per-device mode and clock.}\index{SPI!per-device clock} Because
\texttt{Acquire} sets the mode and clock under the exclusive hold -- a few register
writes while no transfer is in flight -- two devices on one host need not agree on
them. A flash and a display can share \texttt{SPI2}, each clocked its own way:

\begin{lstlisting}
   Setup (SPI2);
   Configure_Pins (SPI2, Sclk => 12, Mosi => 11, Miso => 13);  --  shared wires
   --  ... later, in the flash driver:
   Acquire (S, SPI2, Mode => 0, Clock_Hz =>  8_000_000, CS_Pin => 21);
   --  ... and in the display driver, on the same host:
   Acquire (S, SPI2, Mode => 3, Clock_Hz => 40_000_000, CS_Pin => 10);
\end{lstlisting}

\noindent Each \texttt{Acquire} reprograms the controller for its caller, so the
flash always sees mode 0 at 8\,MHz and the display mode 3 at 40\,MHz -- the SD-card
driver uses exactly this to run its init handshake at 400\,kHz and switch to several
MHz for data. A device wired to \emph{different} pads on the same controller can
likewise pass \texttt{Sclk}/\texttt{Mosi}/\texttt{Miso} to \texttt{Acquire}, which
re-routes the GPIO matrix for the hold; left at \texttt{No\_Pin} they keep the shared
\texttt{Configure\_Pins} routing.

\section{I2C -- master}\index{I2C driver}

\begin{lstlisting}
procedure Setup (Host : I2C_Host; Clock_Hz : Positive := 100_000);
procedure Configure_Pins (Host : I2C_Host; Scl, Sda : Pin_Id);
procedure Acquire (S : in out Session; Host : I2C_Host);
procedure Write (S : Session; Addr : Slave_Address; Data : Byte_Array;
                 Success : out Boolean; Check_Ack : Boolean := True);
procedure Read  (S : Session; Addr : Slave_Address; Data : out Byte_Array;
                 Success : out Boolean);
\end{lstlisting}

\textbf{Parameters.}
\begin{description}[style=nextline]
\item[\texttt{Host}] \texttt{I2C0} or \texttt{I2C1}.
\item[\texttt{Clock\_Hz}] SCL frequency; default \texttt{100\_000} (standard mode),
  \texttt{400\_000} for fast mode.
\item[\texttt{Scl}, \texttt{Sda}] the two bus lines (\texttt{Pin\_Id}); the bus is
  open-drain, so fit external pull-ups.
\item[\texttt{Addr}] 7-bit slave address, \texttt{0\,..\,16\#7F\#}
  (\texttt{Slave\_Address}).
\item[\texttt{Data}] the bytes to write, or the \texttt{out} buffer to read into
  (\texttt{Byte\_Array} -- its length sets the count).
\item[\texttt{Success}] \texttt{out Boolean} -- \texttt{True} if the slave
  acknowledged.
\item[\texttt{Check\_Ack}] (\texttt{Write}) require an ACK on each byte; default
  \texttt{True}.
\end{description}

\textbf{Example 1 -- write a register on a sensor at address \texttt{16\#68\#}.}
\begin{lstlisting}
with ESP32S3.I2C; use ESP32S3.I2C;
procedure I2C_Write_Reg is
   S  : Session;
   Ok : Boolean;
begin
   Setup (I2C0, Clock_Hz => 400_000);
   Configure_Pins (I2C0, Scl => 9, Sda => 8);
   Acquire (S, I2C0);
   Write (S, Addr => 16#68#, Data => (16#6B#, 16#00#), Success => Ok);  --  reg<=0
end I2C_Write_Reg;
\end{lstlisting}

\textbf{Example 2 -- register read (write the address, then read).}
\begin{lstlisting}
with ESP32S3.I2C; use ESP32S3.I2C;
procedure I2C_Read_Reg is
   S    : Session;
   Ok   : Boolean;
   Who  : Byte_Array (0 .. 0);
begin
   Setup (I2C0, Clock_Hz => 400_000);
   Configure_Pins (I2C0, Scl => 9, Sda => 8);
   Acquire (S, I2C0);
   Write (S, 16#68#, (1 => 16#75#), Ok);     --  point at WHO_AM_I register
   Read  (S, 16#68#, Who, Ok);               --  read it back
end I2C_Read_Reg;
\end{lstlisting}

\section{UART -- asynchronous serial}\index{UART driver}

\begin{lstlisting}
--  Startup (single-threaded): the ONLY port-based configuration call.
procedure Setup (Port : UART_Port; Baud : Baud_Rate := 115_200;
                 Bits : Data_Bits := 8; Parity : Parity_Mode := None;
                 Stop : Stop_Bits := One;
                 Tx, Rx, Rts, Cts : Optional_Pin := No_Pin;  --  routes the pads too
                 Rx_Flow_Threshold : Natural := 100);
procedure Acquire (S : in out Session; Port : UART_Port);
procedure Write (S : Session; Data : Byte_Array);
procedure Read  (S : Session; Data : out Byte_Array; Count : out Natural);
function  Available (S : Session) return Natural;
--  Post-Setup reconfiguration -- each takes the HELD Session, not a Port:
procedure Set_Baud      (S : Session; Baud   : Baud_Rate);
procedure Set_Data_Bits (S : Session; Bits   : Data_Bits);
procedure Set_Parity    (S : Session; Parity : Parity_Mode);
procedure Set_Stop_Bits (S : Session; Stop   : Stop_Bits);
procedure Configure_Pins  (S : Session; ...);   --  re-route pads
procedure Enable_Loopback (S : Session; On : Boolean := True);
\end{lstlisting}

\textbf{Parameters.}
\begin{description}[style=nextline]
\item[\texttt{Port}] \texttt{UART0}, \texttt{UART1} or \texttt{UART2}.
\item[\texttt{Baud}] bit rate (\texttt{Baud\_Rate}, \texttt{300\,..\,5\_000\_000});
  default \texttt{115\_200}.
\item[\texttt{Bits}] data bits per frame (\texttt{Data\_Bits}, \texttt{5\,..\,8});
  default \texttt{8}.
\item[\texttt{Parity}] \texttt{None} (default), \texttt{Even} or \texttt{Odd}.
\item[\texttt{Stop}] stop bits: \texttt{One} (default) or \texttt{Two}.
\item[\texttt{Tx}, \texttt{Rx}, \texttt{Rts}, \texttt{Cts}] the four pads as
  \texttt{Optional\_Pin}; all default to \texttt{No\_Pin} (route only the lines the
  link uses).
\item[\texttt{Rx\_Flow\_Threshold}] RX FIFO fill, in bytes, at which RTS
  de-asserts; default \texttt{100} (the FIFO holds 128).
\item[\texttt{On}] (\texttt{Enable\_Loopback}) \texttt{True} (default) folds TX back
  to RX internally.
\item[the \texttt{Set\_*} setters] each take the held \texttt{Session} and one new
  value of the matching \texttt{Setup} type (e.g. \texttt{Set\_Parity} takes a
  \texttt{Parity\_Mode}); \texttt{Configure\_Pins} additionally accepts per-line
  \texttt{*\_Invert} booleans (default \texttt{False}).
\end{description}

\texttt{Setup} both configures the frame format \emph{and} routes the four pins
(all optional -- pass only what the link uses, or none for an internal loopback
self-test). It is the one port-based call, meant to run once at startup before
any task contends for the port. \emph{Every} later change -- baud, data bits,
parity, stop bits, pin re-routing, loopback -- takes a \texttt{Session} instead
of a \texttt{Port}, so you must \texttt{Acquire} the port before you can
reconfigure it. Changing a parameter therefore requires owning the port and can
never race another task's transfer. Each frame setter is an independent
read-modify-write that leaves the other attributes and the pin routing untouched.

This rule is enforced \emph{structurally}, not by convention. Inside the body,
the per-port register handle lives in a nested package whose body hides it; the
only way to obtain it is one ownership-checked gateway that raises
\texttt{Not\_Owned} unless the passed \texttt{Session} holds the port. Every
transfer and every reconfiguration goes through that gateway, so a call without
ownership fails loudly -- and, just as importantly, a \emph{new} operation
added to the driver cannot reach the registers without passing a \texttt{Session}
through the same check. The guard is impossible to forget because there is no
other door to the hardware.

\textbf{Example 1 -- transmit a line.}
\begin{lstlisting}
with ESP32S3.UART; use ESP32S3.UART;
procedure UART_Tx is
   S : Session;
begin
   Setup (UART1, Baud => 115_200, Tx => 17, Rx => 16);   --  format + pins in one
   Acquire (S, UART1);
   Write (S, (Character'Pos ('h'), Character'Pos ('i'), 10));
end UART_Tx;
\end{lstlisting}

\textbf{Example 2 -- echo received bytes (wiring-free, internal loopback).}
\begin{lstlisting}
with ESP32S3.UART; use ESP32S3.UART;
procedure UART_Echo is
   S   : Session;
   Buf : Byte_Array (0 .. 63);
   N   : Natural;
begin
   Setup (UART1);
   Acquire (S, UART1);
   Enable_Loopback (S);              --  TX folded back to RX (held port)
   Write (S, (1 => 16#42#));
   loop
      exit when Available (S) > 0;
   end loop;
   Read (S, Buf, N);                 --  Buf (0 .. N-1) holds the echoed bytes
end UART_Echo;
\end{lstlisting}

\textbf{Example 2, step by step.} This example is worth reading line by line,
because it exercises every part of the \texttt{Session} pattern.

\emph{The declarations (before \texttt{begin}).}
\begin{itemize}
\item \texttt{S : Session} -- \texttt{Session} is a \emph{limited controlled}
  type. \emph{Limited} means it cannot be copied or assigned (two tasks can
  never share one hold); \emph{controlled} means it has a \texttt{Finalize} that
  runs automatically when \texttt{S} leaves scope. It starts inactive -- it
  owns nothing yet.
\item \texttt{Buf : Byte\_Array (0 .. 63)} -- a 64-element array of
  \texttt{Byte} (which is \texttt{new Interfaces.Unsigned\_8}); the destination
  for the bytes we read back.
\item \texttt{N : Natural} -- receives the count of bytes actually read.
\end{itemize}

\emph{The statements, in execution order.}
\begin{enumerate}
\item \texttt{Setup (UART1)} brings UART1 up with the default frame format
  (115200 baud, 8 data bits, no parity, one stop bit -- the defaults on
  \texttt{Setup}). It opens the controller and records that UART1 is configured.
  No pins are routed: we pass no pin arguments to \texttt{Setup} (they all
  default to \texttt{No\_Pin}) because this is a wiring-free self-test.
\item \texttt{Acquire (S, UART1)} takes exclusive ownership of UART1. It first
  checks the init-before-use guard: had UART1 never been \texttt{Setup}, this
  would raise \texttt{Not\_Initialized}. It then blocks on UART1's protected
  guard until the port is free (here, immediately), marks \texttt{S} active, and
  records that it holds UART1. Until \texttt{S} is released no other task can
  acquire UART1 -- and, because every reconfiguration call below takes the held
  \texttt{Session}, no other task can change UART1's settings either.
\item \texttt{Enable\_Loopback (S)} sets the controller's internal TX$\rightarrow$RX
  loopback bit \emph{on the held port}. Everything the transmitter shifts out is
  fed straight back into the receiver \emph{inside the chip}. Because UART is
  push-pull (unlike I2C's open-drain bus) this loopback is faithful and needs
  \emph{no external wiring and no GPIO pads} -- the whole point of the example.
  It takes \texttt{S} (not \texttt{UART1}): you can only flip loopback on a port
  you own.
\item \texttt{Write (S, (1 => 16\#42\#))} -- \texttt{(1 => 16\#42\#)} is a
  one-element array aggregate, a \texttt{Byte\_Array} of length 1 holding the
  byte \texttt{16\#42\#} (66, ASCII \texttt{'B'}). \texttt{Write} pushes it into
  the TX FIFO and returns once it is queued. This runs \emph{outside} the
  protected lock -- the lock only arbitrated \emph{ownership} in the previous
  step; the blocking I/O is never done while holding it. With loopback on, that
  byte will reappear in the RX FIFO shortly.
\item The poll loop \texttt{loop exit when Available (S) > 0; end loop;} spins,
  re-reading how many bytes sit in the RX FIFO, until at least one has arrived.
  The looped-back byte takes roughly one character-time ($\sim$87\,\textmu s at
  115200 baud), so the loop spins briefly then exits. (A self-test idiom;
  production code would use a timeout or an interrupt instead of a bare spin.)
\item \texttt{Read (S, Buf, N)} reads up to \texttt{Buf'Length} (64) bytes into
  \texttt{Buf} and sets \texttt{N} to how many it actually got -- here
  \texttt{N = 1}, with \texttt{Buf (0) = 16\#42\#}. Elements
  \texttt{Buf (N .. 63)} are left untouched. \texttt{Read} waits only briefly
  per byte and does a short read on timeout, so it never hangs on bytes that
  never come.
\item At \texttt{end UART\_Echo}, \texttt{S} goes out of scope. Being controlled,
  its \texttt{Finalize} runs automatically and calls \texttt{Release}, handing
  UART1 back to its guard. This happens \emph{even if an exception had been
  raised} between \texttt{Acquire} and here -- that is why the RAII handle
  exists: the lock cannot leak.
\end{enumerate}

The two phases are deliberately separate. \texttt{Setup} (step 1) is the one
port-based call, meant to run single-threaded at startup before any task
contends for the port. \texttt{Acquire} (step 2) is the per-transaction
ownership gate that hands you a \texttt{Session} -- and everything after it,
including \texttt{Enable\_Loopback} here, acts through that \texttt{Session}, so
a reconfiguration is only possible while you hold the port. Two guards tie the
phases together: you cannot \texttt{Acquire} a port you never \texttt{Setup}, and
you cannot change a port's parameters without holding it.

\textbf{Example 3 -- change baud and frame format independently.} Each frame
attribute has its own setter, so a protocol that renegotiates the link (or a
bootloader that bumps the speed after handshaking) can change just one without
re-running \texttt{Setup} or touching the pin routing. The setters take the held
\texttt{Session}, so you \texttt{Acquire} the port first -- the reconfiguration
runs under the same ownership as the transfers:
\begin{lstlisting}
with ESP32S3.UART; use ESP32S3.UART;
procedure UART_Reconfig is
   S : Session;
begin
   Setup (UART1, Tx => 17, Rx => 16);   --  115200 8-N-1 to start
   Acquire (S, UART1);                  --  own the port before reconfiguring
   Set_Baud      (S, 9_600);            --  then change each attribute alone
   Set_Data_Bits (S, 7);
   Set_Parity    (S, Even);
   Set_Stop_Bits (S, Two);              --  now 9600 7-E-2
   Write (S, (1 => 16#55#));
end UART_Reconfig;
\end{lstlisting}
